\chapter{الضرب}\label{ch:g3:mult}

ثلاثة أكياس في كل واحد سبع كرات: يمكنك أن تجمع $7 + 7 + 7$، لكن
هناك طريقة أسرع --- وهي الضرب، عملية الجمع المتكرر. يبني هذا
الفصل جداول الضرب من الصور، ثم يعلّم طريقة
الأعمدة.

\section{ما معنى الضرب}

\begin{definition}[الضرب]\label{def:g3:mult:def}
\omterm{def:g2:mult:def}{الجداء} $3 \times 7$ يعني «ثلاث مرات سبعة»:
\[
3 \times 7 = 7 + 7 + 7 = 21 .
\]
وفي الصورة، $3 \times 7$ مستطيل من النقط: $3$ صفوف من $7$
نقط.
\end{definition}

\begin{omfigure}
\begin{tikzpicture}[scale=0.42]
  \foreach \r in {0,1,2}
    \foreach \c in {0,...,6}
      \fill[omDef] (\c,-\r) circle (0.27);
  \node[right, font=\small] at (7.0,-1.0) {$3$ صفوف من $7$:\ \
  $3 \times 7 = 21$};
  \begin{scope}[yshift=-4.7cm]
  \foreach \r in {0,...,6}
    \foreach \c in {0,1,2}
      \fill[omThm] (\c,-\r) circle (0.27);
  \node[right, font=\small] at (7.0,-3.0) {$7$ صفوف من $3$:\ \
  $7 \times 3 = 21$ أيضًا!};
  \end{scope}
\end{tikzpicture}

{\small النقط $21$ نفسها، معدودة بالصفوف أو بالأعمدة: فترتيب
\omterm{def:g2:mult:def}{الجداء} لا يهمّ. ومعرفة $3 \times 7$ تعطي
$7 \times 3$ مجانًا --- أي نصف الجداول!}
\end{omfigure}

\begin{proposition}[قواعد تساعد]\label{prop:g3:mult:rules}
\begin{enumerate}
  \item الترتيب لا يهمّ: $a \times b = b \times a$.
  \item والضرب في $1$ لا يغيّر شيئًا؛ والضرب في $0$ يعطي
        $0$.
  \item والضرب في $10$ يُلحق صفرًا: $34 \times 10 = 340$؛
        وفي $100$، صفرين.
  \item ويمكن إعادة بناء جدول من جيرانه:
        $6 \times 8 = 5 \times 8 + 8 = 40 + 8 = 48$.
\end{enumerate}
\end{proposition}
\admitted

\begin{example}[مربّع جداول الضرب]\label{ex:g3:mult:table}
تجتمع الجداول كلها في مربّع واحد. ولقراءة $6 \times 7$: السطر
$6$، والعمود $7$.
\begin{center}
\renewcommand{\arraystretch}{1.15}
\begin{tabular}{c|ccccccccc}
$\times$ & $1$ & $2$ & $3$ & $4$ & $5$ & $6$ & $7$ & $8$ & $9$ \\
\hline
$1$ & 1 & 2 & 3 & 4 & 5 & 6 & 7 & 8 & 9 \\
$2$ & 2 & 4 & 6 & 8 & 10 & 12 & 14 & 16 & 18 \\
$3$ & 3 & 6 & 9 & 12 & 15 & 18 & 21 & 24 & 27 \\
$4$ & 4 & 8 & 12 & 16 & 20 & 24 & 28 & 32 & 36 \\
$5$ & 5 & 10 & 15 & 20 & 25 & 30 & 35 & 40 & 45 \\
$6$ & 6 & 12 & 18 & 24 & 30 & 36 & 42 & 48 & 54 \\
$7$ & 7 & 14 & 21 & 28 & 35 & 42 & 49 & 56 & 63 \\
$8$ & 8 & 16 & 24 & 32 & 40 & 48 & 56 & 64 & 72 \\
$9$ & 9 & 18 & 27 & 36 & 45 & 54 & 63 & 72 & 81 \\
\end{tabular}
\end{center}
والمربّع متناظر حول قطره --- وتلك هي القاعدة 1 مرئيةً
للعيان.
\end{example}

\section{الضرب بالأعمدة}

\begin{method}[الضرب في عدد من رقم واحد]\label{met:g3:mult:column}
لحساب $234 \times 6$:
\begin{enumerate}
  \item اِضرب الآحاد: $4 \times 6 = 24$: اُكتب $4$ واحتفظ برقم $2$؛
  \item واِضرب العشرات وأضف المحتفظ به:
        $3 \times 6 + 2 = 20$: اُكتب $0$ واحتفظ برقم $2$؛
  \item واِضرب المئات وأضف المحتفظ به:
        $2 \times 6 + 2 = 14$: اُكتب $14$.
\end{enumerate}
\[
\begin{array}{r}
2\,3\,4 \\
\times \quad 6 \\
\hline
1\,4\,0\,4
\end{array}
\]
والتقدير للتحقّق: العدد $234$ قريب من $200$، و
$200 \times 6 = 1\,200$: فالجواب $1\,404$ معقول.
\end{method}

\begin{example}[التفكيك للضرب الذهني]\label{ex:g3:mult:split}
طريقة الأعمدة تفكّك العدد سرًّا حسب المراتب. وذهنيًّا، اِفعل
الشيء نفسه:
\[
23 \times 4 = (20 + 3) \times 4 = 20 \times 4 + 3 \times 4
= 80 + 12 = 92 .
\]
ومع جار ودود: $99 \times 5 = 100 \times 5 - 5 =
495$.
\end{example}

\begin{omfigure}
\begin{tikzpicture}[scale=0.55]
  \fill[omDef!18] (0,0) rectangle (10,4);
  \fill[omThm!18] (10,0) rectangle (11.5,4);
  \draw[thick] (0,0) rectangle (11.5,4);
  \draw[thick] (10,0) -- (10,4);
  \node[font=\small] at (5,2) {$20 \times 4 = 80$};
  \node[font=\small, rotate=90] at (10.75,2) {$3 {\times} 4$};
  \node[below, font=\small] at (5,0) {$20$};
  \node[below, font=\small] at (10.75,0) {$3$};
  \node[left, font=\small] at (0,2) {$4$};
\end{tikzpicture}

{\small لماذا ينجح التفكيك: مستطيل $23 \times 4$ من المربعات
الواحدية هو جزء $20 \times 4$ زائد جزء $3 \times 4$ ---
أي $80 + 12 = 92$ مربعًا.}
\end{omfigure}

\section{الضرب في المسائل}

\begin{example}\label{ex:g3:mult:problem}
تحمل حافلة صغيرة $9$ ركاب. فكم راكبًا تحمل $27$
حافلة صغيرة؟
\begin{enumerate}
  \item إنها عملية ضرب: $27$ مجموعة من $9$.
  \item $27 \times 9 = 243$ (بالأعمدة، أو
        $27 \times 9 = 27 \times 10 - 27 = 270 - 27$).
  \item والجملة: تستطيع الحافلات أن تحمل $243$ راكبًا.
\end{enumerate}
\end{example}

\section{تمارين}

\begin{exercise}[$\star$]\label{exo:g3:mult:1}
اُكتب على صورة ضرب، ثم احسب:
$8 + 8 + 8 + 8$؛ \quad $5 + 5 + 5$؛ \quad
$20 + 20 + 20 + 20 + 20$.
\end{exercise}

\begin{exercise}[$\star$]\label{exo:g3:mult:2}
اِسرد وأكمل: $6 \times 7$؛ \ $8 \times 4$؛ \ $9 \times 9$؛
\ $7 \times 8$؛ \ $6 \times 6$.
\end{exercise}

\begin{exercise}[$\star$]\label{exo:g3:mult:3}
اُرسم مستطيل نقط يمثل $4 \times 6$، ثم آخر يمثل $6 \times 4$.
فماذا تلاحظ في العددين؟
\end{exercise}

\begin{exercise}[$\star$]\label{exo:g3:mult:4}
احسب مستعملًا قواعد \cref{prop:g3:mult:rules}:
$56 \times 10$؛ \quad $8 \times 100$؛ \quad $73 \times 0$؛ \quad
$45 \times 1$؛ \quad $30 \times 10$.
\end{exercise}

\begin{exercise}[$\star$]\label{exo:g3:mult:5}
احسب بالأعمدة: $132 \times 3$؛ \quad $217 \times 4$؛ \quad
$408 \times 7$.
\end{exercise}

\begin{exercise}[$\star$]\label{exo:g3:mult:6}
احسب ذهنيًّا بالتفكيك (كما في \cref{ex:g3:mult:split}):
$31 \times 5$؛ \quad $42 \times 3$؛ \quad $25 \times 6$.
\end{exercise}

\begin{exercise}[$\star$]\label{exo:g3:mult:7}
احسب بجار ودود: $99 \times 7$؛ \quad
$101 \times 6$؛ \quad $98 \times 5$.
\end{exercise}

\begin{exercise}[$\star$]\label{exo:g3:mult:8}
قدِّر أولًا، ثم احسب: $389 \times 5$ (قدِّر بالعدد $400$)؛
و $612 \times 8$ (قدِّر بالعدد $600$).
\end{exercise}

\begin{exercise}[$\star$]\label{exo:g3:mult:9}
تسع علبة البيض $6$ بيضات. فكم بيضة في $38$ علبة؟ اُكتب
العملية، واحسب، وأجب بجملة كاملة.
\end{exercise}

\begin{exercise}[$\star\star$]\label{exo:g3:mult:10}
في غرفة $8$ صفوف من $4$ طاولات؛ وتجلس على كل طاولة $2$ من
التلاميذ. فكم تلميذًا تتسع له الغرفة؟ (ضربان.)
\end{exercise}

\begin{exercise}[$\star\star$]\label{exo:g3:mult:11}
جِد كل طرق ترتيب $24$ كرسيًّا في صفوف متساوية (صفوف من
$1$، ومن $2$، \dots). فأي الترتيبات ممكنة؟ (اِستعمل مربّع
الجداول: اِبحث عن $24$ داخله.)
\end{exercise}

\begin{exercise}[$\star\star$]\label{exo:g3:mult:12}
تحسب زينة $47 \times 6$ هكذا: $40 \times 6 = 240$،
و $7 \times 6 = 42$، ثم $240 + 42 = 282$. وتحسب لينا
$47 \times 6 = 50 \times 6 - 3 \times 6 = 300 - 18 = 282$. اِشرح
كل طريقة. وأيّهما تفضّل في $38 \times 4$؟ احسبه
بالطريقتين.
\end{exercise}
